%\SweaveUTF8




\setkeys{Gin}{width=\textwidth}

\begin{Problem}[Planning a stratified sample]\label{problem:stsi_plan}
    We \emph{plan} a survey to estimate the population mean of environmental
    expenditures (\code{expend}). We want to assess whether stratification by
    industry group (\code{code}) can improve precision relative to a simple
    random sample without replacement (\SRS).

    \vspace{0.5em}

    We use the data set \code{sepe.dat} from Problem Set 1. Treat the data as
    an \SRS of $n = 800$ establishments from a population of size $N = 3770$. We
    focus on industry class 31, that is, the domain $U_{31}$ of size $N_{31} =
    650$, identified with \code{code == 1}. Assume that:
    \begin{itemize}
        \item the planned net sample size is $n=800$ establishments;
        \item the anticipated response rate is 80\,\% in both strata; and
        \item the within-stratum variances of \code{empl} and \code{expend} can
            be estimated from the existing \SRS \; \code{sepe.dat}.
    \end{itemize}

\begin{Subtask}
    %---------------------------------------------------------------------------
    \item Allocate the net sample \emph{proportionally} to the two strata.
        Determine the required gross sample sizes and estimate the variance of
        the stratified estimator of the population mean of expenditures.

    \begin{Solution}
        The population sizes are $N_{rest} = 3120$ for \code{code == 0} and
        $N_{31} = 650$ for \code{code = =1}, so $N = 3770$.


\begin{Schunk}
\begin{Sinput}
> sepe <- read.table("http://stat.ethz.ch/Teaching/Datasets/WBL/sepe.dat",
                     header = TRUE)
\end{Sinput}
\end{Schunk}

\begin{Schunk}
\begin{Sinput}
> Nh <- c("rest" = 3120, "31" = 650)
> N <- sum(Nh)
> Wh <- Nh / N
> Wh
\end{Sinput}
\begin{Soutput}
     rest        31 
0.8275862 0.1724138 
\end{Soutput}
\end{Schunk}

        Under proportional allocation,
        \begin{equation*}
            n_h \approx n \, \Wh, \qquad \Wh = \frac{\Nh}{N}.
        \end{equation*}
        We round the stratum sample sizes while keeping their sum equal to
        $n = 800$.

\begin{Schunk}
\begin{Sinput}
> n <- 800
> n_prop <- round(n * Wh)
> n_prop["rest"] <- n - n_prop["31"]
> n_prop
\end{Sinput}
\begin{Soutput}
rest   31 
 662  138 
\end{Soutput}
\end{Schunk}

        Because the anticipated response rate is the same in both strata, the
        gross sample can also be allocated proportionally. The total gross
        sample size is $800/0.8 = 1000$.

\begin{Schunk}
\begin{Sinput}
> response_rate <- 0.8
> n_gross <- round(n / response_rate)
> n_prop_gross <- round(n_gross * Wh)
> n_prop_gross["rest"] <- n_gross - n_prop_gross["31"]
> n_prop_gross
\end{Sinput}
\begin{Soutput}
rest   31 
 828  172 
\end{Soutput}
\end{Schunk}

        Thus, the required gross sample sizes are
        828 establishments with \code{code == 0} and
        172 with \code{code == 1}.

        The variance of the stratified mean $\mystrat = \sum_{h = 1}^H \Wh \,
        \myh$ is
        \begin{equation*}
            \var(\mystrat) = \sum_{h = 1}^H \Wh^2
              \left(1 - \frac{n_h}{\Nh}\right)\frac{\varyh}{n_h}.
        \end{equation*}
        We estimate $\varyh$ using the existing simple random sample.

\begin{Schunk}
\begin{Sinput}
> sigma2_y_h <- tapply(sepe$expend, sepe$code, var)
> var_prop <- sum(Wh^2 * (1 - n_prop / Nh) * sigma2_y_h / n_prop)
> var_prop
\end{Sinput}
\begin{Soutput}
[1] 4423859
\end{Soutput}
\end{Schunk}

        The estimated variance under proportional allocation is therefore
        4'423'859.

        \begin{Remark}
            The response rate affects the \emph{gross} sample size, but the
            variance calculation is based on the planned net sample sizes.
            Using the general variance formula also avoids any approximation
            caused by rounding the proportional allocation.
        \end{Remark}
    \end{Solution}

    %---------------------------------------------------------------------------
    \vspace{0.25em}
    \item Now use a Neyman-type (optimum) allocation, but base the allocation
        on the within-stratum standard deviations of the auxiliary variable
        \code{empl} rather than on \code{expend}.
        \begin{itemize}
            \item Determine the required gross sample sizes.
            \item Under what conditions is it reasonable to derive the
                allocation from an auxiliary variable?
            \item Estimate the variance of the stratified estimator of mean
                environmental expenditures under this allocation.
        \end{itemize}

    \begin{Solution}
        The Neyman allocation for a target variable $y$ is
        \begin{equation*}
            n_h = n \, \frac{\Wh \, \sigma_{y,h}^2}{\sum_{j = 1}^H W_j \,
                \sigma_{y,j}^2}, \qquad h = 1, \ldots, H.
        \end{equation*}

        Here we deliberately replace $\sigma_{y,j}^2$ by the (estimates of)
        within-stratum standard deviations of \code{empl}.

\begin{Schunk}
\begin{Sinput}
> sigma_x_h <- tapply(sepe$empl, sepe$code, sd)
> names(sigma_x_h) <- c("rest", "31")
> alloc_x <- Wh * sigma_x_h
> n_aux <- round(n * alloc_x / sum(alloc_x))
> n_aux
\end{Sinput}
\begin{Soutput}
rest   31 
 629  171 
\end{Soutput}
\begin{Sinput}
> n_aux["rest"] <- n - n_aux["31"]
> n_aux
\end{Sinput}
\begin{Soutput}
rest   31 
 629  171 
\end{Soutput}
\end{Schunk}

        The corresponding gross allocation is

\begin{Schunk}
\begin{Sinput}
> n_aux_gross <- round(n_gross * alloc_x / sum(alloc_x))
> n_aux_gross["rest"] <- n_gross - n_aux_gross["31"]
> n_aux_gross
\end{Sinput}
\begin{Soutput}
rest   31 
 787  213 
\end{Soutput}
\end{Schunk}

        An auxiliary-variable allocation is useful when the pattern of
        within-stratum variability of the auxiliary variable resembles that
        of the target variable. More precisely, we would like $\sigma_{y,h}
        \approx c \cdot \sigma_{x,h}$ with approximately the same positive
        constant $c$ across strata. In that case, the allocation based on $x$
        is close to the Neyman allocation for $y$. Historical values of the
        target variable are a common example of a useful auxiliary variable.

        The existing sample can be used to inspect this condition directly:

\begin{Schunk}
\begin{Sinput}
> sigma_y_h <- sqrt(sigma2_y_h)
> cbind(sigma_x_h = sigma_x_h, sigma_y_h = sigma_y_h,
        ratio = sigma_y_h / sigma_x_h)
\end{Sinput}
\begin{Soutput}
     sigma_x_h sigma_y_h    ratio
rest   177.936  62525.73 351.3945
31     231.377  85396.08 369.0776
\end{Soutput}
\end{Schunk}

        The allocation was computed from \code{empl}, so its variance must be
        evaluated using the \emph{actual} stratum sample sizes above and the
        within-stratum variances of \code{expend}:

\begin{Schunk}
\begin{Sinput}
> var_aux <- sum(Wh^2 * (1 - n_aux / Nh) * sigma2_y_h / n_aux)
> var_aux
\end{Sinput}
\begin{Soutput}
[1] 4332910
\end{Soutput}
\end{Schunk}

        The estimated variance is therefore
        4'332'910.

        % \begin{Remark}
        %     The closed-form minimum variance
        %     \begin{equation*}
        %         \frac{1}{n}\left(\sum_h \Wh \sigma_{y,h}\right)^2
        %         - \frac{1}{N}\sum_h \Wh \sigma_{y,h}^2
        %     \end{equation*}
        %     applies when the allocation itself is based on $\sigma_{y,h}$. It
        %     does not generally apply when the allocation is based on a
        %     different variable such as \code{empl}.
        % \end{Remark}
    \end{Solution}

    %---------------------------------------------------------------------------
    \vspace{0.25em}
    \item Compare the estimated variances from parts a) and b) with the
        estimated variance under a simple random sample of $n=800$. Does
        stratification improve precision?

    \begin{Solution}
        Let $\mys_{\mSRS}$ denote the sample mean for SRS. It holds that
        \begin{equation*}
            \widehat{\var}(\mys_{\mSRS}) =
                \left(1 - \frac{n}{N}\right)\frac{\varyu}{n},
        \end{equation*}
        where $\varyu$ is estimated from \code{sepe.dat}.

\begin{Schunk}
\begin{Sinput}
> var_srs <- (1 - n / N) * var(sepe$expend) / n
> var_srs
\end{Sinput}
\begin{Soutput}
[1] 4567134
\end{Soutput}
\end{Schunk}

        The three planning variances are:

\begin{Schunk}
\begin{Sinput}
> c(SRS = var_srs, proportional = var_prop, auxiliary = var_aux)
\end{Sinput}
\begin{Soutput}
         SRS proportional    auxiliary 
     4567134      4423859      4332910 
\end{Soutput}
\end{Schunk}

        Relative to simple random sampling, the estimated variance reductions
        are

\begin{Schunk}
\begin{Sinput}
> 100 * c(proportional = 1 - var_prop / var_srs,
         auxiliary = 1 - var_aux / var_srs)
\end{Sinput}
\begin{Soutput}
proportional    auxiliary 
    3.137084     5.128461 
\end{Soutput}
\end{Schunk}

        percent. The auxiliary-variable allocation improves on proportional
        allocation if $\widehat{\var}_{\auxalloc} <
        \widehat{\var}_{\propalloc}$. Unlike a true Neyman allocation based on
        \code{expend}, this ordering is not guaranteed because the allocation in
        part b) is based on \code{empl}.
    \end{Solution}
\end{Subtask}
\end{Problem}
